% Time-stamp: <2026-09-05 22:47:07 komori>
% Copyright (c) 2021 Yasushi Komori
% All rights reserved.

%%%%%%%%%%%%%%%%%%%%%%%%%%%%%%%%%%%%%%%%
%
% \brl@bc
% \brl@ebc
% \brl@dbc

\def\brl@bc_appendtoks#1{\brl@bc@toks\expandafter{\the\brl@bc@toks#1}}
\newtoks\brl@bc@toks%

\def\brl@bc_appendextoks#1{\brl@bc@extoks\expandafter{\the\brl@bc@extoks#1}}
\newtoks\brl@bc@extoks%

\def\brl@ebc_appendtoks#1{\brl@ebc@toks\expandafter{\the\brl@ebc@toks#1}}
\newtoks\brl@ebc@toks%

\def\brl@ebc_appendextoks#1{\brl@ebc@extoks\expandafter{\the\brl@ebc@extoks#1}}
\newtoks\brl@ebc@extoks%

\def\brl@dbc_appendtoks#1{\brl@dbc@toks\expandafter{\the\brl@dbc@toks#1}}
\newtoks\brl@dbc@toks

\def\brl@dbc_appendextoks#1{\brl@dbc@extoks\expandafter{\the\brl@dbc@extoks#1}}
\newtoks\brl@dbc@extoks

% default
\def\brl@braille@font{\brl@braille@fontA}% default

%% for pdf
\def\brl@bc_braille#1{%
  {%
    \brl@bc_braille_spread%
    \brl@braille@font%
    \char#1\relax% A or P
  }%
}%

\def\brl@ebc_braille#1#2#3{%
  {%
    \brl@bc_braille_spread%
    \brl@braille@font%
    \char\numexpr#1+#2*64+#3*128\relax%
  }%
}%

% \brl@bc_braille_spread : insert a box when the char is in the position \brl@bse@width.
\def\brl@bc_braille_spread{%
  \ifnum\brl@bse@width=\brl@print@letternum%
  {%
    \brl@orig@color{\brl@char@meancolor}%
    \brl@orig@vrule\@width\brl@braille@width\@height\brl@braille@height\@depth\brl@braille@depth%
  }%
  \fi%
  \global\advance\brl@print@letternum\@ne%
}



%% for nab
\def\brl@bc_bascii{%
  \ifnum\brl@bse@width>\brl@print@letternum%
    \expandafter\brl@bc_bascii_%
  \else%
    \let\brl@bc\brl@bc_bascii_expage%
    \expandafter\brl@bc_bascii_expage%
  \fi%
}%

\def\brl@bc_bascii_#1{%
  \edef\brl@tmp{{\csname brl@bc@bascii_#1\endcsname}}%
  \expandafter\brl@bc_appendtoks\brl@tmp%
  \global\advance\brl@print@letternum\@ne%
}%

\def\brl@bc_bascii_expage#1{%
  \edef\brl@tmp{{\csname brl@bc@bascii_#1\endcsname}}%
  \expandafter\brl@bc_appendextoks\brl@tmp%
  \global\advance\brl@print@letternum\@ne%
}


\def\brl@ebc_ebascii{% extended bascii
  \ifnum\brl@bse@width>\brl@print@letternum%
    \expandafter\brl@ebc_ebascii_%
  \else%
    \let\brl@ebc\brl@ebc_ebascii_expage%
    \expandafter\brl@ebc_ebascii_expage%
  \fi%
}%

\def\brl@ebc_ebascii_#1#2#3{% extended bascii
  % \edef\brl@tmp{{\the\brl@ebc@toks{\the\brl@print@letternum}#2#3}}%
  % \brl@ebc@toks\brl@tmp%
  \edef\brl@tmp{{{\the\brl@print@letternum}#2#3}}%
  \expandafter\brl@ebc_appendtoks\brl@tmp%
  \brl@bc{#1}%
}%

\def\brl@ebc_ebascii_expage#1#2#3{% extended bascii
  % \edef\brl@tmp{{\the\brl@ebc@extoks{\the\numexpr\brl@print@letternum-\brl@bse@width\relax}#2#3}}%
  % \brl@ebc@extoks\brl@tmp%
  \edef\brl@tmp{{{\the\numexpr\brl@print@letternum-\brl@bse@width\relax}#2#3}}%
  \expandafter\brl@ebc_appendextoks\brl@tmp%
  \brl@bc{#1}%
}%

%% for txt (unicode) 
\ifdefined\Uchar% if \Uchar can be used, we use it.
\def\brl@bc_Ucode#1{%
  \expandafter\aftergroup\Uchar\numexpr10240+#1\relax%
}%
\def\brl@ebc_Ucode#1#2#3{%
  \expandafter\aftergroup\Uchar\numexpr10240+#1+#2*64+#3*128\relax%
}%
\else%
\def\brl@bc_Ucode#1{%
  \edef\brl@tmp{\csname brl@bc@Ucode_#1\endcsname}%
  \expandafter\aftergroup\brl@tmp%
}%
\def\brl@ebc_Ucode#1#2#3{%
  \edef\brl@tmp{\csname brl@bc@Ucode_\the\numexpr#1+#2*64+#3*128\relax\endcsname}%
  \expandafter\aftergroup\brl@tmp%
}%
\fi

%% for \DeclareMathOperator
%%
\def\brl@bc_patch#1{%
  \brl@char_appendtoks{\brl@bc{#1}}%
}
\def\brl@ebc_patch#1#2#3{%
  \brl@char_appendtoks{\brl@ebc{#1}#2#3}%
}



%%%%%%%%%%%%%%%%%%%%%%%%%%%%%%%%%%%
%%%%%%%%%%%%%%%%%%%%%%%%%%%%%%%%%%% 
%%
%% \brl@dbc
%%
%% default

%% for pdf
\protected\def\brl@dbc_draw#1{% 
  %
  \brl@bc{0}%
  \brl@@pdf_bcontent%
  \brl@@pdf_invoke{q 1 0 0 1 -\expandafter\brl@@strip@pt\the\brl@braille@width\brl@@space\brl@pgf@dbcht\brl@@space cm}% move \braille@width to the left (back) and translate by \brl@pgf@dbcht
  \brl@dbc_draw_#1\brl@end\brl@end%
}

%% for draw_tikz
\def\brl@dbc_tikz#1{%
  %
  \brl@bc{0}%
  \brl@@pdf_bcontent%
  \brl@@pdf_invoke{q \brl@dbc@scale\brl@@space 0 0 \brl@dbc@scale\brl@@space 0 0 cm}% scale and pt -> bp
  %
  \brl@@pdf_invoke{1 0 0 1 -\expandafter\brl@@strip@pt\the\dimexpr\brl@figure@scale\brl@braille@width\relax\brl@@space\brl@pgf@dbcht\brl@@space cm}% move braille@width to the left (back) and translate by \brl@pgf@dbcht
  \brl@dbc_draw_#1\brl@end\brl@end%
}

\def\brl@dbc_draw_#1#2{%
  \def\brl@tmp{#1}%
  \ifx\brl@tmp\brl@@end%
    \brl@@pdf_invoke{Q}% finished translation
    \brl@@pdf_econtent%
  \else%
    \brl@@pdf_invoke{q 1 0 0 1 #1 -#2 cm}% translation
    \brl@dbc@circle% draw a circle
    \brl@@pdf_invoke{Q}% finish translation
    \expandafter\brl@dbc_draw_%
  \fi%
}


\ifbrl@lualatex%

\newbox\brl@dbc@circle_
\def\brl@dbc_init{%
  \setbox\@ne\hbox{}\wd\@ne\p@\ht\@ne\p@%
  \setbox\z@\hbox{%
    \pdfextension literal{0 g \brl@@cira}%
    \copy\@ne%
  }%
  \saveboxresource type 2 attr{%
/Subtype /Form
/Type /XObject
/FormType 1
/BBox[-1 -1 1 1]}\z@%
  \setbox\brl@dbc@circle_\hbox{\useboxresource\lastsavedboxresourceindex}%
  \wd\brl@dbc@circle_\z@%
  \ht\brl@dbc@circle_\z@%
  \dp\brl@dbc@circle_\z@%
}
\def\brl@dbc@circle{\copy\brl@dbc@circle_}

\else%
  % 
\def\brl@dbc_init{%
  \AtBeginDvi{%
    \special{pdf:bxobj @circle0 bbox -1 -1 1 1}%
    \special{pdf:code 0 g \brl@@cira}%
    \special{pdf:exobj}%
  }
  %
  %
}
\def\brl@dbc@circle{%
  \special{pdf:uxobj @circle0}% draw circle
}%
\fi


%% for print script for ESA721
%% There exists no application which can print out both graphics and texts.
%% The following is a script which gives direct commands to the printer.
%%
%% echo 'commands' > /dev/serialprinter
%%

%% coordinate of ESA721
%%   -> Y
%% |
%% V X
%%
%% for drawing brailles
%% 文字数 行数 X     Y
%% 30     18   0-593 0-449
%% 32     18   0-593 0-479
%% 32     22   0-725 0-479 * use this setting.
%% 40     24   0-792 0-599 ( <- X may be 0-791 (?) )
%%
%% important notices.
%% notice : \dbc@x, \dbc@y employes the ESA721's coordinates
%% notice : after \advance\brl@print@letternum\@ne
%%
\edef\brl@dbc@plotwidth{\the\numexpr\brl@dbc@width*(\brl@bse@width+1)\relax}% 15*(32+1) (the start point of expage)

\let\brl@dbc_appendtmptoks\brl@dbc_appendtoks% default



\def\brl@dbc_output{%
  \ifnum\brl@bse@width>\brl@print@letternum% 32...
    \expandafter\brl@dbc_output_%
  \else% 0...31, i.e., within one page.
    \let\brl@dbc_appendtmptoks\brl@dbc_appendextoks% This change is only local and set the default (appendtoks) after writing one bse line.
    \let\brl@dbc\brl@dbc_output_expage%
    \expandafter\brl@dbc_output_expage%
  \fi%
}

\def\brl@dbc_output_#1{%
  \brl@bc{0}%
  \edef\brl@dbc@y{\the\numexpr\brl@print@letternum*\brl@dbc@width-\brl@dbc@width\relax}% one letter width back.
  \brl@dbc_output_cons#1\brl@end\brl@end%%
}

\def\brl@dbc_output_expage#1{%
  \brl@bc{0}%
  \edef\brl@dbc@y{\the\numexpr\brl@print@letternum*\brl@dbc@width-\brl@dbc@plotwidth\relax}% (one page)+(one letter) width back
  \brl@dbc_output_cons#1\brl@end\brl@end%%
}

\def\brl@dbc_output_cons#1#2{%
  \def\brl@tmp{#1}%
  \ifx\brl@tmp\brl@@end%
  \else%
    {%
      \@tempcnta\numexpr\brl@dbc@y+#1\relax% nonexpr
      \@tempcntb\@tempcnta% Y
      \divide\@tempcnta\@xxxii% [Y/32]
      \@tempcntb\numexpr\@tempcntb-\@tempcnta*\@xxxii+96\relax% Y-[Y/32]*32+96
      \advance\@tempcnta\@xxxii% [Y/32]+32
      \lccode`\!\@tempcnta% YHIGH 0x20-0x3f
      \lccode`\?\@tempcntb% YLOW  0x60-0x7f
      \lowercase{\edef\brl@tmp{\brl@dbc_check_a!\brl@dbc_check_c?}}% YHIGH,YLOW
      %
      \@tempcnta\numexpr\brl@dbc@x+#2\relax% nonexpr
      \@tempcntb\@tempcnta% X
      \divide\@tempcnta\@xxxii% [X/32]
      \@tempcntb\numexpr\@tempcntb-\@tempcnta*\@xxxii+64\relax% X-[X/32]*32+64
      \advance\@tempcnta\@xxxii% [X/32]+32
      \lccode`\!\@tempcnta% XHIGH 0x20-0x3f
      \lccode`\?\@tempcntb% XLOW  0x40-0x5f
      \lowercase{\xdef\brl@tmp{{\brl@tmp\brl@dbc_check_a!\brl@dbc_check_b?}}}% XHIGH,XLOW
      % (0,0) = 0x20 0x60 0x20 0x40 = " ` @"
    }%
    \expandafter\brl@dbc_appendtmptoks\brl@tmp%
    % \show\brl@tmp
    \expandafter\brl@dbc_output_cons%
  \fi%
}

\def\brl@dbc_check_a#1{% 0x20-0x3f
  \if'#1%
    \brl@bse@singleq%
  \else
    #1%
  \fi%
}

{
\catcode`\!=\z@%
!catcode`\\=12!relax%
!gdef!brl@dbc_check_b#1{% 0x40-0x5f
  !if\#1%
    !brl@bse@bs%
  !else%
    #1%
  !fi%
}
}

{
\catcode`\^^?=12\relax%
\gdef\brl@dbc_check_c#1{% 0x60-0x7f
  \if^^?#1%
    \brl@bse@del%
  \else%
    #1%
  \fi%
}
}



\ifbrl@option@pindisplay%
  % Time-stamp: <2024-08-18 20:01:33 komori>
% Copyright (c) 2021 Yasushi Komori
% All rights reserved.

\@namedef{brl@bc_plot@x0}{0}
\@namedef{brl@bc_plot@y0}{2}
\@namedef{brl@bc_plot@x1}{0}
\@namedef{brl@bc_plot@y1}{1}
\@namedef{brl@bc_plot@x2}{0}
\@namedef{brl@bc_plot@y2}{0}
\@namedef{brl@bc_plot@x3}{1}
\@namedef{brl@bc_plot@y3}{2}
\@namedef{brl@bc_plot@x4}{1}
\@namedef{brl@bc_plot@y4}{1}
\@namedef{brl@bc_plot@x5}{1}
\@namedef{brl@bc_plot@y5}{0}


\def\brl@bc_plot#1{%
%  {%
    \count@#1\relax% 
    \ifnum\count@>\z@% if the char is space, skip writing.
      \brl@pgf@dnum\z@%  
      \brl@bc_plot_loop%
    \fi%
    % 
    \global\advance\brl@pgf@ptx\tw@% advance one letter.
    \hbox to 16bp{\brl@orig@vrule\@width\z@\@height\brl@braille@height\hss}%
%  }%
}%


\def\brl@ebc_plot#1#2#3{%
%  {%
    \count@#1\relax%
    \brl@pgf@dnum\z@%  
    \brl@bc_plot_loop%
    % 
    \ifodd#2\relax% 7
      \def\brl@pgf@pointsize{2}%
    \else%
      \def\brl@pgf@pointsize{0}%
    \fi%
    \brl@pgf@currx\brl@pgf@ptx%
%    \brl@pgf@curry\numexpr\brl@pgf@pty-1\relax%
    \brl@pgf@curry\brl@pgf@pty\advance\brl@pgf@curry\m@ne%-1\relax%
    % 
    \brl@pgf_drawpoint%
    % 
    \ifodd#3\relax% 8
      \def\brl@pgf@pointsize{2}%
    \else%
      \def\brl@pgf@pointsize{0}%
    \fi%
    % \brl@pgf@currx\numexpr\brl@pgf@ptx+1\relax%
    % \brl@pgf@curry\numexpr\brl@pgf@pty-1\relax%
    \brl@pgf@currx\brl@pgf@ptx\advance\brl@pgf@currx\@ne%
    \brl@pgf@curry\brl@pgf@pty\advance\brl@pgf@curry\m@ne%
    % 
    \brl@pgf_drawpoint%
    % 
    \global\advance\brl@pgf@ptx\tw@% advance one letter.
    \hbox to 16bp{\brl@orig@vrule\@width\z@\@height\brl@braille@height\hss}%
%  }%
}%


\def\brl@bc_plot_loop{%
  \ifodd\count@%
    \def\brl@pgf@pointsize{2}%
  \else%
    \def\brl@pgf@pointsize{0}%
  \fi%
  % \brl@pgf@currx\numexpr\brl@pgf@ptx+\csname brl@bc_plot@x\the\brl@pgf@dnum\endcsname\relax%
  % \brl@pgf@curry\numexpr\brl@pgf@pty+\csname brl@bc_plot@y\the\brl@pgf@dnum\endcsname\relax%
  \brl@pgf@currx\brl@pgf@ptx\advance\brl@pgf@currx\csname brl@bc_plot@x\the\brl@pgf@dnum\endcsname\relax%
  \brl@pgf@curry\brl@pgf@pty\advance\brl@pgf@curry\csname brl@bc_plot@y\the\brl@pgf@dnum\endcsname\relax%
  \brl@pgf_drawpoint%
  \divide\count@\tw@%
  \advance\brl@pgf@dnum\@ne%
  \ifnum6>\brl@pgf@dnum%
    \expandafter\brl@bc_plot_loop%
  \fi%
}


\endinput


%%% Local Variables:
%%% mode: japanese-latex
%%% backup-each-save-backup : t
%%% TeX-master: t
%%% End:

\else% \ifbrl@option@pindisplay%
  % Time-stamp: <2024-09-26 21:58:28 komori>
% Copyright (c) 2021 Yasushi Komori
% All rights reserved.

\@namedef{brl@bc_plot@x0}{4}
\@namedef{brl@bc_plot@y0}{14}
\@namedef{brl@bc_plot@x1}{4}
\@namedef{brl@bc_plot@y1}{7}
\@namedef{brl@bc_plot@x2}{4}
\@namedef{brl@bc_plot@y2}{0}
\@namedef{brl@bc_plot@x3}{10}
\@namedef{brl@bc_plot@y3}{14}
\@namedef{brl@bc_plot@x4}{10}
\@namedef{brl@bc_plot@y4}{7}
\@namedef{brl@bc_plot@x5}{10}
\@namedef{brl@bc_plot@y5}{0}

%% point: posx, posy
%%
%% 1: 4,14
%% 2: 4,7
%% 3: 4,0
%% 4: 10,14
%% 5: 10,7
%% 6: 10,0
%% 7: 4,-6 4,-11
%% 8: 10,-6 10,-11
%%
%% 15 x 31
%% height : 14
%% depth : 16
%%
%% remove all dots around texts.
%%

\newcount\brl@bc_plot_clear@x
\newcount\brl@bc_plot_clear@y

\def\brl@bc_plot_clear{%
  %
  {%
    \def\brl@pgf@pointsize{0}%
    \brl@bc_plot_clear@x\z@%
    \brl@bc_plot_clear_x%
  }%
}

\def\brl@bc_plot_clear_x{%
  % 0 ... 14 (\brl@plot@braillewidth{15})
%  \brl@pgf@currx\numexpr\brl@pgf@ptx+\brl@bc_plot_clear@x\relax%
  \brl@pgf@currx\brl@pgf@ptx\advance\brl@pgf@currx\brl@bc_plot_clear@x\relax%
  %
  \brl@bc_plot_clear@y\z@%  
  \brl@bc_plot_clear_y%
  \advance\brl@bc_plot_clear@x\@ne%
  \ifnum\brl@plot@braillewidth>\brl@bc_plot_clear@x%
    \expandafter\brl@bc_plot_clear_x%
  \fi%
}

% a letter with 8 dots consists of -11 ... 14
% a letter with 6 dots consists of   0 ... 14
\def\brl@bc_plot_clear_y{%
  % -11 ... -11 + 30 = 19 (\brl@plot@braillewidth{31})
  \brl@pgf@curry\numexpr\brl@pgf@pty-11+\brl@bc_plot_clear@y\relax%
  \brl@pgf_drawpoint%
  % 
  \advance\brl@bc_plot_clear@y\@ne%
  \ifnum\brl@plot@brailleheight>\brl@bc_plot_clear@y%
    \expandafter\brl@bc_plot_clear_y%
  \fi%
}

%% for plotter

\def\brl@bc_plot#1{%
  %
  \ifnum#1>\z@%
    \brl@bc_plot_clear%
  \fi%
  %
  \count@#1\relax%
  \brl@pgf@dnum\z@%  
  \brl@bc_plot_loop%
  %
  \global\advance\brl@pgf@ptx\brl@plot@braillewidth\relax% advance one letter.
  \hbox to \brl@plot@braillewidth bp{\brl@orig@vrule\@width\z@\@height\brl@braille@height\hss}% set the position of the meanings. depth is not needed.
}%


\def\brl@ebc_plot#1#2#3{%
  %
  \ifnum#1>\z@%
    \brl@bc_plot_clear%
  \fi%
  %
  \count@#1\relax%
  \brl@pgf@dnum\z@%  
  \brl@bc_plot_loop%
  %
  \ifodd#2\relax% 7
    % \brl@pgf@currx\numexpr\brl@pgf@ptx+4\relax%
    % \brl@pgf@curry\numexpr\brl@pgf@pty-6\relax%
    \brl@pgf@currx\brl@pgf@ptx\advance\brl@pgf@currx4\relax%
    \brl@pgf@curry\brl@pgf@pty\advance\brl@pgf@curry-6\relax%
    %
    \brl@pgf_drawpoint%
    %
    % \brl@pgf@curry\numexpr\brl@pgf@pty-11\relax%
    \brl@pgf@curry\brl@pgf@pty\advance\brl@pgf@curry-11\relax%
    %
    \brl@pgf_drawpoint%
  \fi%
  \ifodd#3\relax% 8
    %
    % \brl@pgf@currx\numexpr\brl@pgf@ptx+10\relax%
    % \brl@pgf@curry\numexpr\brl@pgf@pty-6\relax%
    \brl@pgf@currx\brl@pgf@ptx\advance\brl@pgf@currx10\relax%
    \brl@pgf@curry\brl@pgf@pty\advance\brl@pgf@curry-6\relax%
    %
    \brl@pgf_drawpoint%
    %
    % \brl@pgf@curry\numexpr\brl@pgf@pty-11\relax%
    \brl@pgf@curry\brl@pgf@pty\advance\brl@pgf@curry-11\relax%
    % 
    \brl@pgf_drawpoint%
  \fi%
  \global\advance\brl@pgf@ptx\brl@plot@braillewidth\relax% advance one letter.
  \hbox to \brl@plot@braillewidth bp{\brl@orig@vrule\@width\z@\@height\brl@braille@height\hss}%
}%

%%
%%

\def\brl@bc_plot_loop{%
  \ifodd\count@%
    % \brl@pgf@currx\numexpr\brl@pgf@ptx+\csname brl@bc_plot@x\the\brl@pgf@dnum\endcsname\relax%
    % \brl@pgf@curry\numexpr\brl@pgf@pty+\csname brl@bc_plot@y\the\brl@pgf@dnum\endcsname\relax%
    \brl@pgf@currx\brl@pgf@ptx\advance\brl@pgf@currx\csname brl@bc_plot@x\the\brl@pgf@dnum\endcsname\relax%
    \brl@pgf@curry\brl@pgf@pty\advance\brl@pgf@curry\csname brl@bc_plot@y\the\brl@pgf@dnum\endcsname\relax%
     %
    \brl@pgf_drawpoint%
    %
  \fi%
  \advance\brl@pgf@dnum\@ne%
  \divide\count@\tw@%
  \ifnum\count@>\z@%
    \expandafter\brl@bc_plot_loop%
  \fi%
}


\def\brl@dbc_plot#1{%
  \brl@bc{0}%
  \def\brl@pgf@pointsize{1}%
  \brl@dbc_plot_#1\brl@end\brl@end%
}

\def\brl@dbc_plot_#1#2{%
  \def\brl@tmp{#1}%
  \ifx\brl@tmp\brl@@end%
  \else%
%    \brl@pgf@currx\numexpr\brl@pgf@ptx-\brl@dbc@width+#1\relax%
    \brl@pgf@currx\numexpr\brl@pgf@ptx-\brl@plot@braillewidth+#1\relax%
    \brl@pgf@curry\numexpr\brl@pgf@pty+\brl@pgf@dbcht-#2\relax%
    % 
    \brl@pgf_drawpoint%
    \expandafter\brl@dbc_plot_%
  \fi%
}



\endinput


%%% Local Variables:
%%% mode: japanese-latex
%%% backup-each-save-backup : t
%%% TeX-master: t
%%% End:

\fi% \ifbrl@option@pindisplay%
%%
%% default
%%
\let\brl@bc\brl@bc_braille%
\let\brl@ebc\brl@ebc_braille%
\let\brl@dbc\brl@dbc_draw%




%%%%%%%%%%%%%%%%%%%%%%%%%%%%%%%%%%%%%%%%

%% 6bit number to braille ascii code
%%
%% \csname brl@bc@bascii_(6bit number)\endcsname = braille ascii code
%%
% UTF+2800 <-> NABCC
%0123456789ABCDEF0123456789ABCDEF0123456789ABCDEF0123456789ABCDEF
% A1B'K2L@CIF/MSP"E3H9O6R^DJG>NTQ,*5<-U8V.%[$+X!&;:4\0Z7(_?W]#Y)=

% 0  1  2  3  4  5  6  7  8  9  A  B  C  D  E  F
% 32 65 49 66 39 75 50 76 64 67 73 70 47 77 83 80
% 34 69 51 72 57 79 54 82 94 68 74 71 62 78 84 81
% 44 42 53 60 45 85 56 86 46 37 91 36 43 88 33 38
% 59 58 52 92 48 90 55 40 95 63 87 93 35 89 41 61 

% "'" (singleeq) if the output is nab, while it is escaped if the output is draw
% "\" (bs) if the output is nab, while it is escaped if the output is draw



{\catcode`\ =11\relax\global\@namedef{brl@bc@bascii_0}{ }}
\@namedef{brl@bc@bascii_1}{A}
\@namedef{brl@bc@bascii_2}{1}
\@namedef{brl@bc@bascii_3}{B}
\@namedef{brl@bc@bascii_4}{\brl@bc@bascii_singleq}% expanded when it is written
\@namedef{brl@bc@bascii_5}{K}
\@namedef{brl@bc@bascii_6}{2}
\@namedef{brl@bc@bascii_7}{L}
\@namedef{brl@bc@bascii_8}{@}
\@namedef{brl@bc@bascii_9}{C}
\@namedef{brl@bc@bascii_10}{I}
\@namedef{brl@bc@bascii_11}{F}
\@namedef{brl@bc@bascii_12}{/}
\@namedef{brl@bc@bascii_13}{M}
\@namedef{brl@bc@bascii_14}{S}
\@namedef{brl@bc@bascii_15}{P}
\@namedef{brl@bc@bascii_16}{"}
\@namedef{brl@bc@bascii_17}{E}
\@namedef{brl@bc@bascii_18}{3}
\@namedef{brl@bc@bascii_19}{H}
\@namedef{brl@bc@bascii_20}{9}
\@namedef{brl@bc@bascii_21}{O}
\@namedef{brl@bc@bascii_22}{6}
\@namedef{brl@bc@bascii_23}{R}
\expandafter\let\csname brl@bc@bascii_24\endcsname\brl@@circum
\@namedef{brl@bc@bascii_25}{D}
\@namedef{brl@bc@bascii_26}{J}
\@namedef{brl@bc@bascii_27}{G}
\@namedef{brl@bc@bascii_28}{>}
\@namedef{brl@bc@bascii_29}{N}
\@namedef{brl@bc@bascii_30}{T}
\@namedef{brl@bc@bascii_31}{Q}
\@namedef{brl@bc@bascii_32}{,}
\@namedef{brl@bc@bascii_33}{*}
\@namedef{brl@bc@bascii_34}{5}
\@namedef{brl@bc@bascii_35}{<}
\@namedef{brl@bc@bascii_36}{-}
\@namedef{brl@bc@bascii_37}{U}
\@namedef{brl@bc@bascii_38}{8}
\@namedef{brl@bc@bascii_39}{V}
\@namedef{brl@bc@bascii_40}{.}
\expandafter\let\csname brl@bc@bascii_41\endcsname\brl@@percent
\@namedef{brl@bc@bascii_42}{[}
\expandafter\let\csname brl@bc@bascii_43\endcsname\brl@@dollar
\@namedef{brl@bc@bascii_44}{+}
\@namedef{brl@bc@bascii_45}{X}
\@namedef{brl@bc@bascii_46}{!}
\expandafter\let\csname brl@bc@bascii_47\endcsname\brl@@ampersand
\@namedef{brl@bc@bascii_48}{;}
\@namedef{brl@bc@bascii_49}{:}
\@namedef{brl@bc@bascii_50}{4}
\@namedef{brl@bc@bascii_51}{\brl@bc@bascii_bs}% expanded when it is written
\@namedef{brl@bc@bascii_52}{0}
\@namedef{brl@bc@bascii_53}{Z}
\@namedef{brl@bc@bascii_54}{7}
\@namedef{brl@bc@bascii_55}{(}
\@namedef{brl@bc@bascii_56}{_}
\@namedef{brl@bc@bascii_57}{?}
\@namedef{brl@bc@bascii_58}{W}
\@namedef{brl@bc@bascii_59}{]}
\expandafter\let\csname brl@bc@bascii_60\endcsname\brl@@sharp
\@namedef{brl@bc@bascii_61}{Y}
\@namedef{brl@bc@bascii_62}{)}
\@namedef{brl@bc@bascii_63}{=}

\endinput


%%% Local Variables:
%%% mode: japanese-latex
%%% backup-each-save-backup : t
%%% TeX-master: t
%%% End:
